\chapter{Escolha de Modelos Gratuitos no OpenCode}
\label{cap:modelos-gratuitos}

\begin{objetivos}
\begin{itemize}
  \item Conhecer os sete modelos gratuitos dispon{\'i}veis no OpenCode e os perfis de cada um;
  \item Selecionar o modelo mais adequado a cada tarefa pedag{\'o}gica com apoio da matriz de decis{\~a}o;
  \item Alternar entre modelos na linha de comando, na interface TUI e na interface gr{\'a}fica.
\end{itemize}
\end{objetivos}

\section{A Democratiza{\c c}{\~a}o do Acesso {\`a} Intelig{\^e}ncia Artificial}

Um dos maiores obst{\'a}culos enfrentados por professores de institui{\c c}{\~o}es p{\'u}blicas de ensino reside no custo financeiro associado {\`a}s assinaturas corporativas de IA e {\`a} cobran{\c c}a por tokens em APIs comerciais pagas em d{\'o}lar.

Para superar essa barreira de acesso e permitir que qualquer docente possa planejar aulas e construir instrumentos avaliativos sem {\^o}nus financeiro, o ecossistema \textbf{OpenCode} disponibiliza uma seleta camada de \textbf{modelos gratuitos da comunidade} (\emph{Free Tier / Community Models}). Esses modelos s{\~a}o mantidos e hospedados por grandes centros de pesquisa, cons{\'o}rcios abertos e pela infraestrutura do pr{\'o}prio ecossistema, oferecendo cotas generosas de requisi{\c c}{\~o}es di{\'a}rias.

Vale registrar que os modelos n{\~a}o s{\~a}o criados pela Anomaly (empresa desenvolvedora do OpenCode): cada um {\'e} propriedade de seu respectivo laborat{\'o}rio de origem -- por exemplo, a fam{\'i}lia \textbf{Nemotron} {\'e} desenvolvida pela \textbf{NVIDIA}, o \textbf{Ling} pelo cons{\'o}rcio aberto \textbf{inclusionAI}, o \textbf{MiMo} pela \textbf{Xiaomi} e o \textbf{Muse Spark} pela \textbf{Meta}. O OpenCode atua apenas como cliente que orquestra o acesso a esses modelos por meio dos identificadores \texttt{opencode/<modelo>}.

Neste cap{\'i}tulo, analisamos detalhadamente os sete modelos gratuitos mais destacados no OpenCode, suas propriedades arquiteturais e as tarefas pedag{\'o}gicas para as quais cada um {\'e} mais indicado. Os subagentes mencionados nas recomenda{\c c}{\~o}es (ex.: \path{@planejador-pedagogico}, \path{@construtor-de-avaliacoes}) s{\~a}o apresentados integralmente no Cap{\'i}tulo~\ref{cap:agentes-skills}.

\section{An{\'a}lise dos 7 Modelos Gratuitos do OpenCode}

\subsection{1. Big Pickle}

\begin{itemize}
  \item \textbf{Identificador}: \path{opencode/big-pickle};
  \item \textbf{Perfil Arquitetural}: Modelo concebido com grande janela de contexto e forte capacidade de s{\'i}ntese enciclop{\'e}dica;
  \item \textbf{For{\c c}as Pedag{\'o}gicas}: Destaca-se na leitura e sumariza{\c c}{\~a}o de longos documentos normativos, projetos pedag{\'o}gicos de curso (PPCs), planos de desenvolvimento institucional (PDIs) e ementas extensas;
  \item \textbf{Quando Usar}: Ideal para a fase inicial de diagn{\'o}stico com o subagente \path{@planejador-pedagogico}, quando o professor precisa cruzar a ementa oficial da disciplina com as compet{\^e}ncias do egresso estabelecidas pelas DCNs (Diretrizes Curriculares Nacionais).
\end{itemize}

\subsection{2. Ling 3.0 Flash Fin Free}

\begin{itemize}
  \item \textbf{Identificador}: \path{opencode/ling-3.0-flash-fin-free};
  \item \textbf{Perfil Arquitetural}: Modelo otimizado para velocidade de infer{\^e}ncia (\emph{flash}) e estrita precis{\~a}o num{\'e}rica, l{\'o}gica e quantitativa;
  \item \textbf{For{\c c}as Pedag{\'o}gicas}: Baix{\'i}ssimo {\'i}ndice de erros aritm{\'e}ticos; excelente manipula{\c c}{\~a}o de tabelas, pontua{\c c}{\~o}es ponderadas e f{\'o}rmulas de avalia{\c c}{\~a}o;
  \item \textbf{Quando Usar}: Perfeito para o c{\'a}lculo e calibra{\c c}{\~a}o de notas em rubricas anal{\'i}ticas complexas, montagem de blueprints de prova com distribui{\c c}{\~a}o de pontos e problemas de engenharia econ{\^o}mica, estat{\'i}stica ou matem{\'a}tica financeira.
\end{itemize}

\subsection{3. Mimo V2.5 Free}

\begin{itemize}
  \item \textbf{Identificador}: \path{opencode/mimo-v2.5-free};
  \item \textbf{Perfil Arquitetural}: Modelo vers{\'a}til, leve e equilibrado, com excelente alinhamento {\`a}s instru{\c c}{\~o}es fornecidas em l{\'i}ngua portuguesa;
  \item \textbf{For{\c c}as Pedag{\'o}gicas}: Reda{\c c}{\~a}o fluida, tom equilibrado (nem rob{\'o}tico, nem informal), facilidade para organizar respostas em t{\'o}picos e alta fidelidade aos esquemas de frontmatter das skills;
  \item \textbf{Quando Usar}: O modelo ``faz-tudo'' para a rotina di{\'a}ria. Recomendado para a reda{\c c}{\~a}o de enunciados de quest{\~o}es discursivas, esbo{\c c}o de cronogramas semanais e media{\c c}{\~a}o de atividades ativas em sala de aula.
\end{itemize}

\subsection{4. Muse Spark 1.2 Free}

\begin{itemize}
  \item \textbf{Identificador}: \path{opencode/muse-spark-1.2-contributor-free};
  \item \textbf{Perfil Arquitetural}: Modelo com calibra{\c c}{\~a}o voltada para idea{\c c}{\~a}o did{\'a}tica, est{\'i}mulo criativo e gera{\c c}{\~a}o de hip{\'o}teses;
  \item \textbf{For{\c c}as Pedag{\'o}gicas}: Capacidade de sugerir analogias pedag{\'o}gicas v{\'a}lidas, met{\'a}foras explicativas e cen{\'a}rios originais do mundo real para problematiza{\c c}{\~a}o contextualizada;
  \item \textbf{Quando Usar}: Na fase de \emph{brainstorming} com a skill \path{ia-educacao-design-problema}, formula{\c c}{\~a}o de situa{\c c}{\~o}es-gatilho para PBL (\emph{Problem-Based Learning}) e din{\^a}micas de debate acad{\^e}mico.
\end{itemize}

\subsection{5. Muse Spark 1.3 Free}

\begin{itemize}
  \item \textbf{Identificador}: \path{opencode/muse-spark-1.3-contributor-free};
  \item \textbf{Perfil Arquitetural}: Evolu{\c c}{\~a}o direta da fam{\'i}lia Muse Spark, combinando a criatividade did{\'a}tica da vers{\~a}o anterior com rigor aprimorado na estrutura de sa{\'i}da;
  \item \textbf{For{\c c}as Pedag{\'o}gicas}: Gera{\c c}{\~a}o de c{\'o}digo \LaTeX{} e tabelas Markdown perfeitamente alinhadas, com consist{\^e}ncia sem{\^a}ntica e rara ocorr{\^e}ncia de alucina{\c c}{\~o}es sint{\'a}ticas;
  \item \textbf{Quando Usar}: Constru{\c c}{\~a}o de rubricas anal{\'i}ticas tabuladas (\path{ia-educacao-rubrica}) e confec{\c c}{\~a}o de exames formais que exigem diagramas ou equa{\c c}{\~o}es matem{\'a}ticas rigorosas.
\end{itemize}

\subsection{6. Nemotron 3 Ultra Free}

\begin{itemize}
  \item \textbf{Identificador}: \path{opencode/nemotron-3-ultra-free};
  \item \textbf{Perfil Arquitetural}: Modelo denso de alt{\'i}ssima capacidade baseado na arquitetura aberta da NVIDIA, com bilh{\~o}es de par{\^a}metros otimizados para racioc{\'i}nio l{\'o}gico avan{\c c}ado;
  \item \textbf{For{\c c}as Pedag{\'o}gicas}: Compreens{\~a}o profunda de c{\'o}digo-fonte (C, C++, Python, Java, Rust), arquitetura de software, teoremas matem{\'a}ticos e demonstra{\c c}{\~o}es formais;
  \item \textbf{Quando Usar}: Avalia{\c c}{\~a}o de projetos t{\'e}cnicos em Ci{\^e}ncia da Computa{\c c}{\~a}o e Engenharias, resolu{\c c}{\~a}o anal{\'i}tica de exerc{\'i}cios de f{\'i}sica avan{\c c}ada e auditoria cr{\'i}tica de gabaritos com o subagente \path{@construtor-de-avaliacoes}.
\end{itemize}

\subsection{7. Nemotron 3.5 Lightning Free}

\begin{itemize}
  \item \textbf{Identificador}: \path{opencode/nemotron-3.5-lightning-free};
  \item \textbf{Perfil Arquitetural}: Variante de resposta quase instant{\^a}nea (\emph{lightning}), balanceando a for{\c c}a l{\'o}gica da arquitetura Nemotron com lat{\^e}ncia ultrabaixa;
  \item \textbf{For{\c c}as Pedag{\'o}gicas}: Devolutivas r{\'a}pidas em fra{\c c}{\~o}es de segundo, mantendo coer{\^e}ncia em disciplinas STEM;
  \item \textbf{Quando Usar}: Sess{\~o}es interativas em tempo real (como a gera{\c c}{\~a}o dinâmica de testes conceituais durante a aula para \emph{Peer Instruction} via \path{ia-educacao-peer-instruction}) e devolu{\c c}{\~a}o r{\'a}pida de feedback para m{\'u}ltiplos estudantes.
\end{itemize}

\begin{itemize}
  \item \textbf{Perfil Arquitetural}: Variante de resposta quase instant{\^a}nea (\emph{lightning}), balanceando a for{\c c}a l{\'o}gica da arquitetura Nemotron com lat{\^e}ncia ultrabaixa;
  \item \textbf{For{\c c}as Pedag{\'o}gicas}: Devolutivas r{\'a}pidas em fra{\c c}{\~o}es de segundo, mantendo coer{\^e}ncia em disciplinas STEM;
  \item \textbf{Quando Usar}: Sess{\~o}es interativas em tempo real (como a gera{\c c}{\~a}o dinâmica de testes conceituais durante a aula para \emph{Peer Instruction} via \path{ia-educacao-peer-instruction}) e devolu{\c c}{\~a}o r{\'a}pida de feedback para m{\'u}ltiplos estudantes.
\end{itemize}

\section{Matriz de Decis{\~a}o Docente: Qual Modelo Escolher?}

A Tabela~\ref{tab:modelos-matriz} resume a recomenda{\c c}{\~a}o de uso de cada modelo gratuito conforme a necessidade pedag{\'o}gica imediata do professor.

\begin{table}[htbp]
  \centering
  \small
  \caption{Matriz de Sele{\c c}{\~a}o de Modelos Gratuitos do OpenCode}
  \label{tab:modelos-matriz}
  \begin{tabularx}{\linewidth}{p{4.8cm}p{4.2cm}X}
    \toprule
    \textbf{Tarefa Pedag{\'o}gica} & \textbf{Modelo Recomendado} & \textbf{Identificador no OpenCode} \\
    \midrule
    An{\'a}lise de PPCs e Ementas Longas & Big Pickle & \path{opencode/big-pickle} \\
    \addlinespace
    C{\'a}lculo de Notas e Pesos & Ling 3.0 Flash Fin Free & \path{opencode/ling-3.0-flash-fin-free} \\
    \addlinespace
    Reda{\c c}{\~a}o de Aulas e Exerc{\'i}cios & Mimo V2.5 Free & \path{opencode/mimo-v2.5-free} \\
    \addlinespace
    Problemas e Estudos de Caso & Muse Spark 1.2 Free & \path{opencode/muse-spark-1.2-contributor-free} \\
    \addlinespace
    Rubricas e Formata{\c c}{\~a}o \LaTeX{} & Muse Spark 1.3 Free & \path{opencode/muse-spark-1.3-contributor-free} \\
    \addlinespace
    Demonstra{\c c}{\~o}es e C{\'o}digo STEM & Nemotron 3 Ultra Free & \path{opencode/nemotron-3-ultra-free} \\
    \addlinespace
    Sess{\~o}es em Sala e Feedback R{\'a}pido & Nemotron 3.5 Lightning Free & \path{opencode/nemotron-3.5-lightning-free} \\
    \bottomrule
  \end{tabularx}
\end{table}

\section{Como Alternar entre Modelos no OpenCode}

O professor pode selecionar ou alternar o modelo em uso de tr{\^e}s formas simples:

\begin{enumerate}
  \item \textbf{Ao Iniciar o OpenCode}: Passando o identificador pela flag \texttt{-m}:
  \begin{lstlisting}[language=bash]
  opencode -m opencode/nemotron-3-ultra-free
  \end{lstlisting}
  \item \textbf{Na Interface Interativa de Terminal (TUI)}: Pressionando a tecla de atalho \texttt{Ctrl+M} (ou digitando \texttt{/model}) e selecionando o modelo na lista;
  \item \textbf{Na Interface Web/Desktop}: Clicando no menu suspenso no canto superior direito da janela e escolhendo o modelo gratuito correspondente com o mouse.
\end{enumerate}
